% Copyright (c) 2026 Twin Vector Labs LLC.
% All rights reserved.
\documentclass[11pt]{article}

\usepackage[T1]{fontenc}
\usepackage{lmodern}
\usepackage[margin=1in]{geometry}
\usepackage{amsmath,amssymb,mathtools}
\usepackage{booktabs}
\usepackage{array}
\usepackage{enumitem}
\usepackage{xcolor}
\usepackage{microtype}
\usepackage[colorlinks=true,linkcolor=blue!50!black,urlcolor=blue!60!black]{hyperref}
\usepackage{xurl}

\newcommand{\C}{\mathbb{C}}
\newcommand{\Z}{\mathbb{Z}}
\newcommand{\rep}[1]{\mathbf{#1}}
\newcommand{\abs}[1]{\lvert #1 \rvert}
\newcommand{\norm}[1]{\lVert #1 \rVert}
\newcommand{\rU}{r_{U}}
\newcommand{\gradS}{\norm{\nabla S_{\mathrm{Regge}}}^2}
% breakable verbatim-ish code/path formatting (xurl: may break at any character)
\DeclareUrlCommand\code{\urlstyle{tt}}
\newcommand{\caveat}[1]{\par\medskip\noindent\textcolor{red!60!black}{\textbf{Caveat.}}\ #1\par\medskip}

\title{\bfseries The fully emergent proton:\\
drive with the residual, observe everything else\\
\large an epic design note, for discussion}
\author{Tessera --- cobordism subsystem\\
(epic proposal; builds on \#457/\#475, \#489/\#491/\#549, \#541)}
\date{}

\begin{document}
\maketitle

\begin{abstract}
\noindent We propose an epic that finishes what the emergent-merge engine started: build the
proton by extremizing the one objective we already have,
$F = \gradS + \Gamma\,\rU$, where $\rU$ pins \emph{only the prepared initial state}
(the experiment's boundary data), and treat \textbf{every other physical quantity ---
color content, flavor, spin, electric charge, mass, radius, even the intermediate event
structure --- as an observable read off the resulting lattice} after convergence. The
current source of truth, \code{Proton::build()}, is most of the way there: topology grows
from a single $\Delta^4$ seed, geometry emerges from the relaxation, nothing is welded or
hand-placed. What remains imposed is a short list --- the two-step event graph with its
colored intermediate targets, the singlet output target, the answer-shaped convergence
gate, and the all-spacelike seed --- and the record says each item can be removed outright
or demoted to a read. The strongest argument is empirical, not aesthetic: hand-imposed
symmetry \emph{provably deleted} the flavor observable (\#414), and the emergent register
\emph{delivered} it (\#475). The design is four independently falsifiable rungs --- the
joint node, the inputs-only objective, Lorentzian activation, and the entangled readout
battery --- each with pre-registered pass/fail criteria whose negatives are results.
\end{abstract}

\section{The goal and the objective}

The objective, stated once (the full \emph{complex} Sorkin/dual Regge action --- keep
$\operatorname{Im}$; extremize $\delta S = 0$, never minimize $\abs{S}$):
\begin{equation}
F \;=\; \gradS \;+\; \Gamma \cdot \rU ,
\qquad
\gradS = \sum_e \Bigl\lvert \frac{\partial S}{\partial \ell_e^2} \Bigr\rvert^2 ,
\label{eq:objective}
\end{equation}
with $\rU$ the relabeling-invariant, zero-filled \code{residualForPeriods} over the
emergent register holes. The epic's thesis: Eq.~\eqref{eq:objective} with $\rU$ restricted
to the \emph{prepared inputs} is the entire drive; the proton is whatever stable structure
emerges, and everything we want to know about it is a measurement on the relaxed lattice.

This note builds directly on the \code{MultiCobordism} engine (\#491/\#549), the T5
emergent-optimizer architecture and findings (\#462/\#475), the emergent-intermediates
bracket (\#434/\#438), the spin/readout theory
(\code{docs/theory/cobordism/proton-spin/}), and the Lorentzian-dormancy investigation
(\#541).

\section{The constitution: what is allowed to shape the lattice}
\label{sec:constitution}

Unchanged from the T5 design note, restated as this epic's ground rules. The \emph{only}
conditions imposed on any build are:
\begin{enumerate}[leftmargin=2em]
  \item the \textbf{manifold gate} --- every move passes \code{dualComplexValid}
        (surgery is allowed \emph{because} it is gated);
  \item \textbf{extremization of $F$} --- greedy $\Delta F$ moves plus the trap door,
        then the Stage-2 Wirtinger relaxation; the action stays complex end to end;
  \item the \textbf{prepared initial state}, held representable through its $\rU$ term
        (never by freezing vertices --- \code{MultiCobordism::pinnedBoundaryVertices} is
        deliberately empty, \code{MultiCobordism.cpp:292});
  \item \textbf{room to act} --- the single-$\Delta^4$ seed plus the gated \code{precone}
        growth knob.
\end{enumerate}
Everything else --- a target topology, a $b_k$ goal, a move recipe, a frame condition,
and (the subject of this epic) a \emph{final-state target} --- is either forbidden or
scheduled for demotion to a post-hoc observable. Every observable claimed on a converged
structure must pass the GAUGE and RELABEL invariance gates (the \#485 convention).

\section[What Proton.cpp still imposes, precisely]{What \texttt{Proton.cpp} still imposes, precisely}
\label{sec:imposed}

\begin{center}
\small
\begin{tabular}{@{}c p{5.6cm} p{3.6cm} p{4.6cm}@{}}
\toprule
\# & imposition & where & tension with the goal \\
\midrule
1 & The \textbf{two-step event graph} with colored intermediate targets: step A drives
    2 neutral pairs $\to$ diquark $\{1,\omega\}$ $\sqcup$ antidiquark $\{1,\omega^2\}$
  & \code{Proton.cpp:67-92}, \code{:143-150}
  & the event structure and a colored intermediate are placed in the objective, not
    observed \\
\addlinespace
2 & The \textbf{hand-off between the steps is nominal}: step B starts a \emph{fresh}
    $\Delta^4$ seed and imposes the \emph{ideal} diquark $+$ third quark $\{\omega^2\}$
    as its inputs; step A's emergent output is never read out and fed forward (only its
    $\rU$ is reported)
  & \code{Proton.cpp:94-114}, \code{:150}
  & as a physical event chain this is not composition (the result-state\,$=$\,next-boundary
    rule is not exercised) \\
\addlinespace
3 & The \textbf{proton singlet $[1,\omega,\omega^2]$ as the output target}, read on the
    whole cobordism
  & \code{Proton.cpp:107-109}; whole-read branch \code{MultiCobordism.cpp:266-273}
  & the answer sits in the objective --- held by $\rU$ (the sanctioned knob), but a drive
    nonetheless \\
\addlinespace
4 & The \textbf{convergence gate} \code{colorResidual < colorTolerance} $\wedge$
    \code{holes} $\geq$ \code{minQuarkHoles} (defaults $0.5$ / $3$)
  & \code{Proton.cpp:159-160}, \code{Proton.h:86-89}
  & selection by the known answer instead of by stationarity $+$ persistence \\
\addlinespace
5 & The \textbf{all-spacelike seed} ($\ell^2 = +1$ everywhere)
  & \code{Proton.cpp:62-63}
  & not a target, but it silently pins the dynamics Euclidean (\S\ref{sec:evidence}.5) ---
    half the observable battery then reads zero for structural reasons \\
\bottomrule
\end{tabular}
\end{center}

Impositions that are \textbf{legitimate and stay}: the prepared inputs (neutral
$\Sigma = 0$ $q\bar q$ pairs --- the choice of experiment), the objective itself, the
gate, and the seed-size knobs.

\section{The evidence for de-imposition}
\label{sec:evidence}

\paragraph{4.1 De-imposition is what unlocked flavor.}
On the hand-symmetric $A_4$ register, flavor is \emph{provably} unreadable: any
relabeling-invariant per-window measurable commutes with the transitive window-cycling
$g$ and is constant on the orbit --- measured discriminator margin $\sim 3\times10^{-6}$,
machine zero (\code{proton_flavor_split.md}, \#414). On the \textbf{emergent,
non-symmetric} register the same read succeeds: per-hole Dirac--K\"ahler charge
distinguishable with relative spread $0.52$--$0.61$, persisting (often sharpening)
through Stage-2 relaxation in $20/21$ structures
(\code{t5_emergent_optimizer_findings.md}, \#475). The lesson generalizes:
hand-imposed symmetry deletes observables; emergence carries them.

\paragraph{4.2 The dynamics does not choose the staged diquark path.}
With both endpoints pinned and the interior free, the emergent intermediate is nearly
neutral (singlet overlap $0.995$, colored $\sigma \approx 0.096$); the imposed
strong-$\bar{\rep{3}}$ diquark is \emph{hosted} (connected-bulk carry $\rU = 0.137$,
$\gradS = 49.1 < 100$) but \emph{diverges} from the emergent path --- overlap $0.519$
against the pre-registered $0.90$
(\code{438-fixed-bipartite-sequence-f5ad527.md}).
\caveat{That bracket ran on the retired $2{+}1$\,D base; the numbers may not transfer.
The emergent-vs-imposed methodology does.}

\paragraph{4.3 The diquark is recoverable as a read, so it need not be a target.}
The pair-loops are Poincar\'e duals of the complementary holes --- $[\gamma_{ij}] = -[k]$
--- so the diquark \emph{is} the dual of the spectator-quark hole, and the
$u{:}u{:}d = 2{:}1$ multiplicity is a dual-basis statement on the finished register
(\code{pair_loop_quarks.tex}, with a concrete falsifiable test: the odd-one-out loop
must coincide with the diquark loop, under the GAUGE/RELABEL gates). Removing the staged
intermediate from the drive loses nothing that cannot be read afterwards.

\paragraph{4.4 The engine already supports every shape this epic needs.}
$N$ input blocks (the ctor takes lists); the \#489 joint build (3 neutral pairs $\to$
[proton, antiproton] as two localized output blocks) ran on this engine and delivered the
carried singlet \emph{and} independent per-hole flavor; and with \textbf{\code{outputTargets}
empty the objective is inputs-only today} --- the \code{rU} else-branch adds no output
term when no output blocks exist (\code{MultiCobordism.cpp:274-284}). Rung 2
(\S\ref{sec:design}) is mechanically a constructor argument, not an engine change.

\paragraph{4.5 The Lorentzian sector is dormant by dynamics, not by physics.}
The converged proton has $\operatorname{Im}\ell^2 \equiv 0$,
$\min\operatorname{Re}\ell^2 = 0.262$ --- the Stage-2 clamp ($[0.05, 20]$,
\code{MultiCobordism.cpp:672-676}) never fires; the all-spacelike seed simply keeps the
descent in the Euclidean basin (\code{541-why-geometry-spacelike-30047e0.md}). Meanwhile
every electric observable is measured to be Lorentzian-gated: emergent Gauss-law charge
$\abs{Q_e} = 0.297\,/\,0.117$ with timelike worldlines vs \textbf{exactly $0$} on the
all-spacelike control (\#434/\#438), and the closed-surface flux is topologically
protected ($\sim 10^{-16}$) either way. Charge, $E/B$, photons-as-null-edges, temporal
curvature, and the transition-amplitude program all wait on this.

\paragraph{4.6 The remaining spin gap is a readout, not a build defect.}
The build creates the entanglement; the per-hole extraction discards it. The product of
per-hole spinors \emph{provably} floors at the uncorrelated mixture ($J^2 \sim 9/4$; the
entangled proton value $3/4$ is unreachable by any Bloch-vector reduction), while the
instrument itself is validated exact on clean states ($3/4$ / $7/4$ / $15/4$) ---
\code{spin_readout.tex}, \code{joint_proton_spin_findings.md} (\#489). The fix is
the total-space operator $J_a = \sum_h P_h \Sigma_a P_h$ whose \textbf{cross-hole terms}
carry the entanglement (\#477), on the fiber$\leftrightarrow$cells lift (\#495);
\code{cartan_weyl_gluon.tex} supplies the gauge-invariant chamber coordinate (the
connected correlator $C_{ij}$) and four pre-scoped experiments, including the decisive
``$C_{ij}$ two ways'' test and the holonomy-sourced chamber point (the operational gluon
test).

\paragraph{4.7 What ``carried vs floored'' already gives us for free.}
Neutral pairs carry ($\rU \sim 10^{-8}$, $\gradS$ at its ${\sim}88$ floor) where
free/colored quarks floor ($\rU \sim 0.5$, $\gradS \sim 14600$) --- confinement as a
property of the objective, established in the merge track. Observation mode leans on
exactly this asymmetry: the drive never says ``singlet'', but the landscape does.

\section{The design: four rungs, each independently falsifiable}
\label{sec:design}

\subsection{Rung 1 --- collapse the two-step into one joint node (the \#489 shape)}
\label{sec:rung1}

One \code{MultiCobordism}: inputs $=$ \textbf{three neutral $q\bar q$ pairs}, outputs
$=$ \textbf{two localized blocks} (baryon $\sqcup$ conjugate antibaryon --- the
multi-output branch of \code{rU}, \code{MultiCobordism.cpp:274-278}), on a single
$\Delta^4$ seed, driven by the same Grow $\to$ directed probes $\to$ Evolve $\to$ Relax
sequence \code{build()} uses today. The two-step build is \textbf{demoted to the
reference oracle} (the \#454 pattern), kept for A-vs-B comparison. The intermediates
become reads: the pair-loop dual-basis flavor test (\S\ref{sec:evidence}.3) and the
mid-build register content replace the imposed diquark target.

\emph{Why two output blocks and not singlet-on-the-whole:} with three pairs the whole
complex must also host the conjugate sector; a single whole-read target
$[1,\omega,\omega^2]$ would fight the antibaryon's $[1,\bar\omega,\bar\omega^2]$
periods. The $2\to2$ recombination branch already models exactly this.

\emph{Discussion item:} if the pairwise-microcausality argument in \code{Proton.h} is to
be preserved as a physical claim, the alternative is to make the two-step's composition
real --- feed step A's \textbf{measured} emergent output forward as step B's boundary
instead of the ideal constants. Both arms can run; the observables adjudicate.

\subsection{Rung 2 --- retire the output targets; convergence $=$ stationarity $+$ persistence}
\label{sec:rung2}

Same node, \textbf{\code{outputTargets} $= \{\}$} --- the objective is
$\gradS + \Gamma \sum_i w\,\rU(\text{input}_i)$, nothing else. The convergence gate
(\S\ref{sec:imposed}, item 4) is replaced by:
\begin{itemize}[leftmargin=2em]
  \item \textbf{stationarity} --- Stage 2 stops on its relative-tolerance test
        (\code{lastStage2Stationary()}), and no Stage-1 candidate move lowers $F$;
  \item \textbf{persistence} --- the claimed structure survives continued Evolve/Relax
        passes and fresh-seed re-relaxations with its observables stable (the Stage-2
        stability-sweep methodology from \#475, promoted to the acceptance criterion).
\end{itemize}
Then \textbf{observe} (\S\ref{sec:battery}): hole count, singlet content \emph{as a
diagnostic} (\code{residualOfTargetStateAgainstHarmonic} already exists as a static ---
used after the fact, it is a measurement, not a drive), $\sigma_R$, per-hole DK charge
spread, spin, charge, mass, radius.

The claim under test, pre-registered: \emph{the stable endpoint of the inputs-only drive
is a bound, localized, color-neutral $\geq 3$-hole object with a conjugate partner --- a
baryon pair --- rather than three identity-transported pairs.} This is genuinely open,
and the record names the failure modes honestly: single-ended pinning under-constrained
the isolated creation node (conformal drift, color spread \emph{worsening}
$0.47 \to 2.44$, \#435), and merge transport tends to identity (\#376). Both outcomes are
results (\S\ref{sec:criteria}).

\subsection{Rung L (parallel) --- activate the Lorentzian sector}
\label{sec:rungL}

Per \#541's recommendation: (a) replace the hard $\operatorname{Re}\ell^2 \geq 0.05$
floor with a \textbf{causal-aware degeneracy guard} --- forbid only the light-cone band
$\abs{\ell^2} < \varepsilon$, both signs, so timelike edges are admissible and degeneracy
still is not; (b) put the time direction in the \textbf{boundary data} --- timelike
worldline edges on the prepared pair inputs (a cobordism has a time direction;
preparation, not answer-insertion). Verify the complex deficit / Lorentzian volume path
across signature changes, and watch the conformal runaway by its restoring force (the
standing rule: never sidestep it). Null edges that appear are the photon candidates, and
they should be \emph{sourced}, not spontaneous.

Unlocks, with their already-built readers: $Q = \oint_S E$ (\code{gaussLawCharge} --- the
genuine holonomy, protected to $\sim 10^{-16}$, vs the drifting hand-weighted covector),
\code{fieldStrengthSplit} ($E/B$), the CPT pairing $Q_p + Q_{\bar p} = 0$ across the
emergent conjugate sectors, the temporal curvature panel, and the action-coupled
amplitude Experiment B of \code{transition_amplitudes_from_relaxation.tex} (does
$\delta S = 0$ protect unitarity under backreaction --- either answer is a result).

\subsection{Rung R (parallel) --- finish the entangled readout battery}
\label{sec:rungR}

Pure reads, no new drive terms, each gated GAUGE/RELABEL:
\begin{enumerate}[leftmargin=2em]
  \item \textbf{\#495 fiber$\leftrightarrow$cells lift} (Whitney/K\"ahler--Atiyah) ---
        the last rung between the validated instrument and a $J^2$ read off the emergent
        proton;
  \item \textbf{\#477 total-space $J^2$} with the cross-hole terms (the entanglement
        carriers);
  \item \textbf{$C_{ij}$ chamber coordinate} $+$ the two cheap
        \code{cartan_weyl_gluon} experiments (instrument reproduction;
        $C_{ij}$-two-ways --- vertical vs holonomy routes, decisive either way);
  \item \textbf{holonomy-sourced chamber point} --- the transport intertwiner $M$ as the
        entangling source; the relaxed proton lands on the proton corner without the
        chamber point being put in by hand (the operational meaning of ``the exponential
        acts like a gluon'');
  \item \textbf{pair-loop dual-basis flavor} on the built proton ($2{:}1$ multiplicity;
        odd-one-out $=$ diquark loop);
  \item \textbf{mass/radius on the relaxed interior}, the \#451 methodology ported to
        4D: closed-fan hinges only, skeleton built in C++, dual-volume radius $+$ shell
        deficit mass, $r \cdot m$ reported with its definitional spread
        (order-of-magnitude claim only).
\end{enumerate}

\subsection{Search practicality (all rungs)}
\label{sec:search}

Removing targets flattens the $\Delta F$ landscape --- T5's cautionary tale: the first
cone-out often raises $F$, and an undirected walk never left $b_2 = 0$ in ${\sim}13.9$\,M
gated moves. The levers are in-tree: the trap door, restarts, the incremental
$\Delta F$, the directed probes, and the RL harness. Two consistency notes.
\code{directedConeOut} scores this node's \code{rU}, which in rung 2 contains only the
input terms --- expect weaker guidance; that is honest, and it is what the policy is for.
And the RL reward for the emergent arm must be the engine's $F$ \textbf{alone} --- the
current \code{rl_objective_search.py} shaping (\code{hole_reward_weight},
\code{rstate_reward_weight}) is fine for the oracle arm but would smuggle the proton
into the policy on the emergent arm. RL work stays on the examples-not-tests lane.

\section{The observable battery (read after convergence; never loop conditions)}
\label{sec:battery}

\begin{center}
\small
\begin{tabular}{@{}p{4.2cm} p{7.2cm} p{3.4cm}@{}}
\toprule
quantity & reader & status \\
\midrule
color content / confinement & \code{residualOfTargetStateAgainstHarmonic} as diagnostic;
  $\sigma_R$; carried-vs-floored contrast & ready \\
flavor & per-hole DK charge spread on the emergent register; pair-loop dual-basis test
  & ready / rung R.5 \\
spin (baryon content) & per-hole $\abs{\langle S\rangle} = \tfrac12$, product $J^2$
  floor $3/2$ (dual-edge frame) & ready \\
spin ($\tfrac12$ vs $\tfrac32$) & total-space $J^2$ with cross-hole terms; $C_{ij}$
  chamber coordinate & rung R.1--R.3 \\
electric charge & $Q = \oint_S E$ Gauss-law holonomy; CPT pairing across conjugate
  sectors & blocked on rung L \\
photon & null-edge census (sourced, not spontaneous) & blocked on rung L \\
mass / radius & interior shell deficit; dual-volume radius; $r \cdot m$
  (order-of-magnitude) & rung R.6 \\
event structure / intermediates & mid-build register content; pair-loop duality
  & rung 1 \\
transition amplitude / unitarity & $Z_T(q,\alpha) = (T^{-1}q)^\dagger G\,\alpha$,
  $u' = (T^\dagger)^{-1}$, $\{G = I,\ T^\dagger T = I\}$ certificates
  & Exp.~A ready; B on rung L \\
\bottomrule
\end{tabular}
\end{center}

\section{Pre-registered criteria, and what each negative would mean}
\label{sec:criteria}

Thresholds set before running, never loosened; negatives are reported as findings (the
\#434/\#438 discipline).

\begin{description}[leftmargin=2em]
  \item[Rung 1.] The joint node converges at least as often as the two-step across a
    fixed seed battery; carried singlet on the baryon block; independent per-hole flavor
    (spread $\geq 0.1$); pair-loop test passes. \emph{Negative $\to$} the two-step's
    staging is load-bearing; keep it and make its composition real
    (\S\ref{sec:rung1}, discussion item).
  \item[Rung 2.] Stationary $+$ persistent endpoint with $\geq 3$ emergent holes on the
    baryon cluster, singlet diagnostic $\leq$ the rung-1 carried level, conjugate pairing
    observed. \emph{Negative (identity transport / no binding) $\to$} the final-state pin
    is demonstrably the minimal necessary imposition --- quantified, and the
    singlet-as-knob (today's \code{Proton.cpp}) stands as the documented fallback.
    \emph{Negative (conformal drift) $\to$} the restoring-force diagnosis from
    \#435/\#434 applies; report which term failed to regulate, do not bolt on a boundary.
  \item[Rung L.] Timelike content survives relaxation (no clamp artifacts --- the guard,
    not the floor, binds); $\abs{Q_e} > 0.05$ with the all-spacelike control at $0$;
    flux protection and CPT pairing hold. \emph{Negative $\to$} the Euclidean basin is
    deeper than the seeding; report the basin boundary, revisit boundary data (not the
    objective).
  \item[Rung R.] Instrument checks reproduce $3/4$ / $7/4$ / $15/4$ exactly; then the
    emergent read either moves $J^2$ decisively below $3/2$ toward $3/4$ (proton
    channel) or lands at $15/4$ ($\Delta$) --- \emph{either definite channel is a
    success}; an indefinite ${\sim}9/4$ after the cross-hole terms would falsify the
    total-space-operator hypothesis itself, which is worth knowing.
  \item[Control (suggested).] Run rung 2 with \textbf{two} pairs --- does a meson-shaped
    2-hole neutral object emerge? Distinguishes ``binding is generic under $F$'' from
    ``three is special'' (the bipartite obstruction says the singlet is irreducibly
    3-valent; a meson control probes the complementary claim).
\end{description}

\section{Out of scope}

\begin{itemize}[leftmargin=2em]
  \item The codim-2 surgical move / $b_2$ register (proven unreachable and mooted ---
        flavor reads at $k = 3$).
  \item Any auxiliary term or rule in $F$ ($b_k$ goals, harmonic-dimension policing,
        frame conditions).
  \item Retiring the hole-based period readout cores or other load-bearing API
        refactors.
  \item Docs/animation work beyond the findings notes each rung owes.
\end{itemize}

\section{Proposed epic decomposition (tickets cut after sign-off)}
\label{sec:tickets}

\begin{center}
\small
\begin{tabular}{@{}l l p{11.6cm}@{}}
\toprule
ticket & rung & deliverable \\
\midrule
E1  & 1 & joint formation node as a \code{Proton} build mode driving the \#489 shape;
          two-step demoted to reference oracle; A-vs-B harness $+$ findings note \\
E2  & 1 & pair-loop dual-basis flavor read on the built proton (GAUGE/RELABEL-gated) \\
E3  & 2 & inputs-only build (\code{outputTargets} $= \{\}$); stationarity $+$
          persistence acceptance; observable-battery report; meson control \\
E4  & L & causal-aware degeneracy guard replacing the $\operatorname{Re}\ell^2$ floor;
          Lorentzian boundary data on the pair worldlines; signature-change
          verification \\
E5  & L & charge battery on the Lorentzian build: $Q = \oint E$, $E/B$, null-edge
          census, CPT pairing across conjugate sectors \\
E6  & R & \#495 fiber$\leftrightarrow$cells lift $+$ \#477 total-space $J^2$ (C++ on
          \code{MultiCobordism}/\code{EigenstateSynthesis}) \\
E7  & R & $C_{ij}$ chamber readout: instrument check $+$ $C_{ij}$-two-ways $+$ emergent
          read \\
E8  & R & holonomy-sourced chamber point (the gluon test) \\
E9  & R & mass/radius battery on the relaxed 4D interior (\#451 methodology) \\
E10 & --- & RL emergent arm: reward $=$ engine $F$ only; benchmark vs random/grow-only
          on E3 (examples lane) \\
E11 & L & transition-amplitude Experiment A on the equivariant fill, then B on the
          relaxed build (unitarity under backreaction) \\
\bottomrule
\end{tabular}
\end{center}

Dependencies: E1 $\to$ E2/E3; E4 $\to$ E5/E11-B; E6 $\to$ E7 $\to$ E8; E9, E10 ride on
E1/E3. Suggested order of first bites: E1 $+$ E4 in parallel (E1 is mostly
configuration; E4 is the physics unlock), E2 immediately after E1 (cheap, pure read).

\section{Open questions for the discussion}
\label{sec:questions}

\begin{enumerate}[leftmargin=2em]
  \item \textbf{Two-step disposition} --- oracle only, or also make its composition real
    (feed the measured step-A output forward)? Keeping both arms is cheap; the header's
    physical claim (\code{Proton.h:22-33}) deserves an explicit verdict rather than
    silent retirement.
  \item \textbf{Input specification for the joint node} --- the $\Z_3$-symmetric neutral
    triple ($\{1,-1,0\}$, $\{0,1,-1\}$, $\{-1,0,1\}$) vs the current pair set; and how
    much \code{inputResidualWeight} / $\Gamma$ re-tuning rung 2 needs once the output
    terms vanish from \code{rU}'s scale.
  \item \textbf{Persistence protocol} --- how many continued passes / fresh-seed
    re-relaxations, and which observables must be stable at what tolerance, for
    ``persistent'' to be claimed (candidate: the \#475 sweep protocol verbatim).
  \item \textbf{Lorentzian boundary data} --- timelike $\ell^2$ on which edges exactly
    (the pair worldlines only, or the whole seed given a CDT-style time direction), and
    the guard band $\varepsilon$.
  \item \textbf{Claim language} --- at what point do we say ``proton'' rather than
    ``baryon-shaped bound state''? Proposal: ``proton'' requires the rung-R channel read
    ($J^2 \to 3/4$) \emph{and} the charge read ($+1$ vs the neutron's $0$), both of
    which are observations, not targets.
\end{enumerate}

\section*{References}

\begin{itemize}[leftmargin=2em]
  \item \code{src/cobordism/Proton.cpp}, \code{include/cobordism/Proton.h} --- the
        current two-step source of truth.
  \item \code{src/cobordism/MultiCobordism.cpp} (\code{rU} :257,
        \code{pinnedBoundaryVertices} :292, Stage-2 clamp :672),
        \code{include/cobordism/MultiCobordism.h} --- the engine.
  \item \code{docs/design/t5_emergent_optimizer_design.md},
        \code{t5_emergent_optimizer_findings.md} --- the architecture and the
        emergent-flavor result.
  \item \code{docs/design/434-emergent-intermediates-f1fb86d.md},
        \code{438-fixed-bipartite-sequence-f5ad527.md}, \code{435-bipartite-creation.md}
        --- the emergent-vs-imposed bracket and the under-constraint lesson.
  \item \code{docs/design/541-why-geometry-spacelike-30047e0.md} --- the Lorentzian
        dormancy finding and fix.
  \item \code{docs/design/proton_flavor_split.md}, \code{proton_charge_gauss_law.md},
        \code{proton_electroweak_finding.md}, \code{proton_radius_finding.md},
        \code{451-proton-geometry-0a1a5aa.md} --- the observable readers and their
        obstructions.
  \item \code{docs/design/composite_proton_spin_findings.md},
        \code{joint_proton_spin_findings.md},
        \code{docs/theory/cobordism/proton-spin/spin_readout.tex},
        \code{pair_loop_quarks.tex}, \code{cartan_weyl_gluon.tex} --- the spin
        program.
  \item \code{docs/design/transition_amplitudes_from_relaxation.tex} --- the
        amplitude/unitarity layer.
  \item \code{docs/design/proton_bipartite_obstruction.tex},
        \code{proton_symmetric_windows.tex}, \code{proton_session_review.tex} ---
        the representation-theoretic ground ($\epsilon_{ijk}$ is irreducibly 3-valent)
        and the symmetric-window era.
  \item \code{examples/cobordism/rl_objective_search.py} --- the RL harness whose
        reward the emergent arm restricts to $F$.
\end{itemize}

\end{document}
